\documentclass[12pt]{article}
\usepackage{float}
\usepackage{graphicx}
\usepackage{fontspec}
\usepackage{url}
\usepackage{apacite}
\setmainfont{Times New Roman}
\usepackage[letterpaper,top=2cm,bottom=2cm,left=2cm,right=2cm,marginparwidth=1.5cm]{geometry}

\title{Weekly Assignment 9: Evaluating methods for enchasing cybersecurity behaviour related to different types of insiders in organizations}
\author{Truong Khang Tobias Nguyen}
\date{}

\begin{document}

\maketitle

\section{Prioritized Belief Groups}
An effective cybersecurity environment does not only depend on technical solutions, policies, but also on how an organization's employees perceive security measures. Employee's beliefs have an influence on their behaviour on security-related issues. Negatively perceived actions may lead to non-compliant behaviour, which in turn turn can lead to further security measures and subsequently undermining employee's attitude even further \cite{balozian2023adversarial}.
\newline
That being said, while I believe that all groups need to be addressed, special attention needs to be given to groups whose behaviour can negatively impact an organization's security efforts. That being said, I believe that organizations should focus particularly on \textbf{cynical} and \textbf{sceptical} groups. 
\newline
Cynical employees may consider security measures unnecessary which in turn can result in them disregarding measures as they are apathetic to them \cite{malik2026security}. On the other hand, sceptical employees may question the value of security measures and may resist measures they regard as ineffective or pointless. Both groups require attention because increasing enforcement may strenghten their negative beliefs, become non-compliant or encourage workarounds \cite{balozian2023adversarial}.
\newline
While I stated that organizations should put special attention to these two groups, acquiescent and convinced groups should not be neglected. While acquiescent employees may comply with security measures, they do it simply because they are required to rather than through inherent motivation or understanding. Although this behaviour may seem secure at first, compliance through obligation becomes fragile, especially when security measures conflict with work objectives.
\newline
Convinced employees, however, understand the impact of security incidents and subsequent necessity of appropriate measures and are thus more likely to follow security measures. In that sense, although they do not require special attention by organizations, organization should maintain their engagement and actively utilize them as security advocates and accelerators.

\section{Methods, Interventions and Theories}
In order to sway \textbf{cynical} groups towards secure behaviour, organizations should replace generic warnings and measures with tangible, specific examples of how security incidents affect their organization. According to the Protection Motivation Theory, individuals are more likely to adopt secure behaviour when they recognize threats, risks and believe that recommended measures are effective \cite{zis-Rogers1983Cognitive}. In that sense, showing the value of security measures may directly challenge their assumptions. However, organizations should also investigate employees’ frustrations, why they have them and address them accordingly. Such an approach may help understanding the beliefs and experiences that this group has and possibly alleviate them.
\newline
\textbf{Acquiescent} groups need to understand why security measures exist and the possible impact on the organization if they did not have them. Mandatory training and measures creates superficial compliance without actual commitment or understanding. Organizations should instead provide interactive training, realistic examples, and practical demonstrations explaining why certain security measures are in place. The Theory of Planned Behavior argues that  attitudes, perceived  norms and behavioural control are essential in shaping intentions \cite{ajzen1991theory}. Training that builds employees’ understanding in security can help transform rule-following into informed and consistent behaviour.
\newline
As \textbf{sceptical} understand the importance of security without believing in their measures, organizations should actively involve them in the process of developing security measures and explicitly invite them to participate. Employee's perception of the legitimacy of security measures influences compliance and thus, it is vital that sceptical groups are involved \cite{cram2023legitimacy}. Explaining the reasoning for measures, acknowledging criticisms, and involving employees can likely be more effective than relying on sanctioning non-compliant behaviour because it addresses the reasons employees question security measures.
\newline
Lastly, as \textbf{confident} groups are already aware of the impact of security incidents and their required measures, organizations should strenghten their commitment and enable them. Recognizing and acknowledging secure behaviour, provide advanced training and enabling them as security champions are methods to maintain their motivation and engagement. These employees can develop good practices, support colleagues and subsequently positively influence other parts of the organization. To support this, the Theory of Planned Behavior suggests that supportive social norms can strengthen intentions to behave securely \cite{ajzen1991theory}. 

\section{Conclusion}
To conclude this report, although organizations should have different approaches for the individual groups, these approaches should be inclusive. Cynical employees need credible evidence; acquiescent employees need understanding; sceptical employees need participation; and convinced employees need reinforcement and opportunities to influence others. Across all groups, security measures must be usable, proportionate, and supported by each group. 

\bibliographystyle{apacite}
\bibliography{references}

\end{document}